\documentclass[a4paper,11pt]{article}
\usepackage[utf8]{inputenc}
\usepackage[T1]{fontenc}
\usepackage[spanish,es-noshorthands]{babel}
\usepackage[margin=2.5cm]{geometry}
\usepackage{booktabs}
\usepackage{parskip}
\usepackage{enumitem}
\usepackage{xcolor}
\definecolor{vino}{RGB}{90,21,21}
\usepackage{sectsty}
\allsectionsfont{\color{vino}}

\title{\color{vino}\textbf{Cat\'alogo Online --- Informe de avance}\\[4pt]
\large Santa Mar\'ia Distribuidora \--- Los Polvorines}
\author{}
\date{6 de julio de 2026}

\begin{document}
\maketitle
\vspace{-2.5em}

\section*{Resumen ejecutivo}

El cat\'alogo online (\texttt{don-florencio.vercel.app}) ya funciona como \textbf{vidriera digital con venta por WhatsApp}: los clientes navegan los productos con precios minorista y mayorista, y consultan o piden por WhatsApp con un toque. El sitio est\'a \textbf{t\'ecnicamente listo}; lo que falta para difundirlo son sobre todo \textbf{decisiones comerciales} (validar precios, dominio propio) que se detallan al final.

\section{Qu\'e se hizo}

\subsection*{1. Reorganizaci\'on completa del cat\'alogo}
\begin{itemize}[leftmargin=1.2em, itemsep=2pt]
  \item Se cargaron e integraron \textbf{1.107 productos} (incluyendo los 600+ de Almac\'en).
  \item \textbf{Todos} los nombres fueron prolijados (antes ven\'ian del sistema de facturaci\'on, en may\'usculas y con c\'odigos) y \textbf{cada producto tiene ahora una descripci\'on}.
  \item Se orden\'o todo en \textbf{14 categor\'ias} claras: Fiambres, Quesos, Embutidos, L\'acteos, Almac\'en, Bebidas, Congelados, Conservas, Limpieza, \textbf{Papelera} (renombrada a pedido de la direcci\'on), Panader\'ia, Pastas, Aderezos y Dulces.
  \item Se eliminaron 23 productos duplicados y se corrigieron errores puntuales de nombre y categor\'ia.
\end{itemize}

\subsection*{2. Precios}
\begin{itemize}[leftmargin=1.2em, itemsep=2pt]
  \item Todos los productos tienen \textbf{precio minorista}.
  \item Se carg\'o el \textbf{precio mayorista} calculado como \textbf{minorista con 12\% de descuento} (criterio provisorio, ver ``Decisiones'').
  \item El cliente puede alternar entre ver precios minoristas o mayoristas en el sitio.
\end{itemize}

\subsection*{3. Fotos de productos}
\begin{itemize}[leftmargin=1.2em, itemsep=2pt]
  \item \textbf{601 de 1.107 productos (54\%) ya tienen foto} con fondo blanco, estilo supermercado.
  \item Se consiguieron autom\'aticamente m\'as de 200 fotos oficiales de marcas conocidas (Arcor, La Seren\'isima, Knorr, La Virginia, etc.), todas revisadas una por una antes de publicar.
  \item Las restantes son de marcas locales sin foto en internet: se cargan a mano desde el panel de administraci\'on (cada producto tiene un link de b\'usqueda que facilita el trabajo).
\end{itemize}

\subsection*{4. Velocidad, seguridad y calidad}
\begin{itemize}[leftmargin=1.2em, itemsep=2pt]
  \item La carga del cat\'alogo baj\'o de \textbf{27 segundos a $\sim$1 segundo}.
  \item Se hicieron \textbf{dos auditor\'ias completas} del sistema; se corrigieron todos los problemas detectados (incluido un error por el cual ``eliminar producto'' no borraba de verdad).
  \item El sitio funciona como \textbf{app instalable} en el celular y qued\'o protegido con las medidas de seguridad web est\'andar. Hay \textbf{copias de respaldo} de todos los datos.
\end{itemize}

\subsection*{5. Marca y contacto}
\begin{itemize}[leftmargin=1.2em, itemsep=2pt]
  \item WhatsApp actualizado a \textbf{+54 9 11 5723-8769} en todo el sitio.
  \item \textbf{Instagram} conectado (\texttt{@santamaria\_dist}).
  \item Horarios unificados: \textbf{Lun--Vie 7 a 21 h, S\'ab 7 a 13 h}.
  \item Al compartir el link por WhatsApp ahora se ve la vista previa con el logo.
\end{itemize}

\section{N\'umeros clave}

\begin{table}[h]
\centering
\begin{tabular}{lr}
\toprule
\textbf{Indicador} & \textbf{Valor} \\
\midrule
Productos publicados & 1.107 \\
Categor\'ias & 14 \\
Productos con foto & 601 (54\%) \\
Productos con descripci\'on & 100\% \\
Precios minorista / mayorista & 100\% / 100\% \\
Tiempo de carga del cat\'alogo & $\sim$1 segundo \\
\bottomrule
\end{tabular}
\end{table}

\section{Decisiones que necesitamos de la direcci\'on}

\begin{enumerate}[leftmargin=1.4em, itemsep=4pt]
  \item \textbf{Margen mayorista:} hoy es un 12\% de descuento sobre el minorista, igual para todo. ?`Est\'a bien as\'i o var\'ia por categor\'ia (ej.\ fiambres vs.\ almac\'en)? Se recalcula en minutos.
  \item \textbf{Precios minoristas:} vienen del sistema de facturaci\'on. ?`Est\'an vigentes o hay que actualizar antes de difundir?
  \item \textbf{Domingos:} el sitio dice que se abre de lunes a s\'abado. ?`Se abre los domingos? (Antes figuraba ``Dom 8--12''.)
  \item \textbf{Dominio propio:} hoy la direcci\'on es \texttt{don-florencio.vercel.app}. Se recomienda comprar un dominio propio (ej.\ \texttt{santamariadistribuidora.com.ar}, costo bajo anual) para tarjetas, Instagram y confianza del cliente.
  \item \textbf{Pr\'oximos cambios:} quedamos a la espera de los dem\'as ajustes que la direcci\'on quiera conversar.
\end{enumerate}

\section{Plan sugerido para salir a la venta}

\begin{enumerate}[leftmargin=1.4em, itemsep=2pt]
  \item Validar precios (minorista vigente + margen mayorista real). \hfill \textit{direcci\'on}
  \item Revisar stock: que lo publicado refleje lo que hay. \hfill \textit{direcci\'on + panel}
  \item Comprar y conectar el dominio propio. \hfill \textit{1 d\'ia}
  \item Crear el perfil de \textbf{Google Business} (aparecer en Maps y b\'usquedas locales, gratis). \hfill \textit{1 d\'ia}
  \item Difundir: Instagram, estados de WhatsApp, tarjetas con el QR del sitio.
  \item En paralelo: seguir completando fotos de las categor\'ias m\'as vistas.
\end{enumerate}

\vspace{1em}
\begin{center}
\color{vino}\rule{0.6\textwidth}{0.4pt}\\[6pt]
\small Sitio: \texttt{don-florencio.vercel.app} \quad---\quad WhatsApp: +54 9 11 5723-8769 \quad---\quad IG: @santamaria\_dist
\end{center}

\end{document}
